% HD-F1-T16. GPT-6 Astra Pro, /root/transport; 22 November 2026 UTC.
\section[HD-F1-T16: Lie-type defining fields below 137]{HD-F1-T16: every remaining almost simple Lie-type source has defining field size below 137}
\label{sec:f1-large-field-lie-source-exclusion}

In the remaining T2-OTHER natural-curve source gate, an almost
simple projective coefficient image whose socle is of Lie type
has defining field size \(q<137\). This assertion concerns the
original coefficient representation over each of the eight
remaining imaginary quadratic fields. It includes arbitrary
central lifts and all inner and outer inertia elements. There
is no irreducibility assumption on restriction to the perfect
central layer of the socle.

More precisely, an actual almost simple projective source of
this kind with \(q\ge137\), parabolic dimension six and period
rank nine would have original coefficient rank \(m\le25\).
The remaining gate requires \(m\ge27\), so no such source
occurs. The defining field parameter is the usual \(q=p^a\)
in \({}^2B_2(q),{}^2G_2(q),{}^2F_4(q)\) as well as in
the other Lie-type families.

\paragraph{1. The precise published character input and its central-cover scope.}
Use M.~W. Liebeck, \emph{Character ratios for finite groups
of Lie type, and applications}, \emph{Contemp. Math.}~694
(2017), 193--208. In the notation of its Section~2,
\(\widetilde S=\mathcal S^F\) is a quasisimple fixed-point
group of a simply connected simple algebraic group. The
proof of Proposition~2.7, published p.~198 (author PDF p.~6),
states the following combined consequence of its Theorem~2.4:
\[
 \frac{|\chi(x)|}{\chi(1)}<\frac3{\sqrt q}
 \quad
 \bigl(x\in\widetilde S\setminus Z(\widetilde S),\quad
       1\ne\chi\in\operatorname{Irr}_{\mathbf C}(\widetilde S)\bigr).
 \tag{F1T16.1}
\]
The combined bound is the input used here. The source record
distinguishes the inspected published survey passage from
the original Gluck proof, whose full text was unavailable
in this run.

For \(q\ge137\), this simply connected fixed-point group
is the universal central cover of its simple central
quotient \(S\). The exact classification input is
Malle--Testerman, \emph{Linear Algebraic Groups and Finite
Groups of Lie Type} (2011), Theorem~24.17 and
Remark~24.19 with Table~24.3, pp.~213--214.
The exceptions to perfectness have defining field size
at most three, and every entry in the exceptional-cover
table has defining field size at most nine. The complete
exceptional-multiplier list is also printed in
Hiss--Husen--Magaard, \emph{Imprimitive irreducible modules
for finite quasisimple groups}, \emph{Mem. Amer. Math.
Soc.}~234 (2015), no.~1104, Chapter~5, Table~1, p.~35.
Only these stated classification inputs, rather than an
identification of arbitrary central extensions, are used
below.

\paragraph{2. The actual central layer and a surjective comparison map.}
Let \(G\subset GL_K(V)\) be the original faithful
absolutely irreducible finite coefficient image, put
\(m=\dim_KV\), and let \(Z=Z(G)\), its scalar centre.
The given almost simple projective image is
\(G/Z\subset\operatorname{Aut}(S)\). Let \(P\) be the
inverse image of its simple socle \(S\), and put \(L=P'\).
Since \(S\) is perfect, \(P'\) surjects onto \(S\), so
\[
 P=ZL,\qquad
 L=[P,P]=[ZL,ZL]=[L,L],\qquad
 Z(L)=L\cap Z.
 \tag{F1T16.2}
\]
For the last equality, a central element of \(L\) maps
to the centre of \(S\), which is trivial; the converse
holds because \(Z\) is scalar. Thus \(L\to S\) is the
actual perfect central extension arising inside \(G\).
It is normal in \(G\). The subspace
\(V_{\mathbf C}^{L}\) is consequently \(G\)-stable.
Absolute irreducibility, together with faithfulness and
the nontrivial quotient \(L\to S\), gives
\[
                         V_{\mathbf C}^{L}=0.
 \tag{F1T16.3}
\]

The universal property supplied by paragraph~1 gives an
actual homomorphism
\[
 \beta:\widetilde S\longrightarrow L,
 \qquad
 \bigl(\widetilde S\xrightarrow{\beta}L\longrightarrow S\bigr)
       =\bigl(\widetilde S\longrightarrow S\bigr).
 \tag{F1T16.4}
\]
It is surjective. Indeed its image \(H\) maps onto \(S\),
so \(L=Z(L)H\); perfectness then gives
\(L=[L,L]=[H,H]\subseteq H\). Its kernel is central,
because it lies in the kernel of the map to \(S\).
This is a coefficient-representation comparison map.
No lift of the original curve cover to a
\(\widetilde S\)-cover is asserted or needed.

\paragraph{3. The full restricted character, including multiplicities.}
Write the actual semisimple restriction as
\[
 V_{\mathbf C}|_L=\bigoplus_i a_i W_i,
 \qquad a_i\in\mathbf Z_{>0},\qquad
 m=\sum_i a_i\dim W_i.
 \tag{F1T16.5}
\]
By (F1T16.3), every irreducible \(W_i\) is nontrivial.
The surjection \(\beta\) makes its inflation to
\(\widetilde S\) irreducible and nontrivial. For any
\(h\in L\setminus Z(L)\), a preimage under \(\beta\)
is noncentral, because its image in \(S\) is nontrivial.
Apply (F1T16.1) to these inflations and add their actual
characters with their multiplicities. With
\(r_q=3/\sqrt q\), one obtains
\[
 |\operatorname{Tr}_{V}(h)|
 \le\sum_i a_i|\operatorname{Tr}_{W_i}(h)|
 <r_q\sum_i a_i\dim W_i=r_qm.
 \tag{F1T16.6}
\]
Finite order makes \(h\) diagonalizable with eigenvalues
of absolute value one. For
\(c_h=\operatorname{rank}(h-1)\), the fixed eigenvalues
contribute \(m-c_h\) and the other eigenvalues have
sum of absolute value at most \(c_h\). Therefore
\[
 m-2c_h\le|\operatorname{Tr}_{V}(h)|<r_qm,
 \qquad
 \operatorname{rank}(h-1)>\frac{1-r_q}{2}m.
 \tag{F1T16.7}
\]

\paragraph{4. Every actual outer lift and every eigenspace.}
Let \(g\in G\) be nonscalar, and let \(x\ne1\) be
its image in \(G/Z\subset\operatorname{Aut}(S)\).
The centralizer of \(S\) in \(\operatorname{Aut}(S)\)
is trivial: an automorphism commuting with every inner
automorphism sends each \(s\in S\) to an element with
the same inner automorphism, hence to \(s\), because
\(Z(S)=1\). Choose \(a\in S\) with \([x,a]\ne1\)
and an actual lift \(e\in L\). Normality of \(L\)
and this nontrivial projective commutator show that
\[
               h=[g,e]=geg^{-1}e^{-1}
                       \in L\setminus Z(L).
 \tag{F1T16.8}
\]
For every \(\lambda\in\mathbf C\), the exact matrix
identity is
\[
 h-1=(ge-eg)g^{-1}e^{-1}
     =\bigl((g-\lambda)e-e(g-\lambda)\bigr)g^{-1}e^{-1}.
 \tag{F1T16.9}
\]
The rank of this actual map is at most
\(2\operatorname{rank}(g-\lambda)\). Consequently
\[
 \operatorname{rank}(g-\lambda)
       >\frac{1-3/\sqrt q}{4}m
 \quad\text{for every nonscalar }g\in G
       \text{ and every }\lambda\in\mathbf C.
 \tag{F1T16.10}
\]
This includes elements inducing field, graph or diagonal
automorphisms, all orders and all their scalar lifts.
A nonidentity scalar \(g\) has
\(\operatorname{rank}(g-1)=m\), which satisfies the same
strict moving-rank bound for \(q\ge137\).

\paragraph{5. The original branch budget gives the numerical contradiction.}
For the original connected full sphere cover, let
\(g_1,\ldots,g_b\) be its true inertia matrices. The
original parabolic dimension and period dimension give
\[
 \sum_{j=1}^b\operatorname{rank}(g_j-1)=2m+6,
                         \qquad b\ge12.
 \tag{F1T16.11}
\]
Applying (F1T16.10) and the scalar case to these same
matrices yields
\[
 2m+6>
 \frac{b(1-3/\sqrt q)}4m
 \ge3(1-3/\sqrt q)m,
 \qquad
 m<\frac6{1-9/\sqrt q}.
 \tag{F1T16.12}
\]
For \(q\ge137\) the denominator is positive and
\[
 \frac9{\sqrt q}\le\frac9{\sqrt{137}}<\frac{10}{13},
 \qquad
 117^2=13689<13700=100\cdot137.
 \tag{F1T16.13}
\]
It follows that \(m<26\), hence \(m\le25\). This
contradicts \(m\ge27\) in the remaining source gate
and proves the stated exclusion.

\paragraph{6. Original source, primitive corner and framework maps.}
The group homomorphism \(\beta\) is used only to inflate
characters and estimate the actual matrices in
\(V_{\mathbf C}\). Equations (F1T16.8)--(F1T16.9)
are commutator and matrix maps in that original space.
The map of covers is the identity of
\(C\to\mathbf P^1\), and the map of middle-extension
sheaves is its identity, with the same branch stalks
\(V^{g_j}\). Its Euler characteristic gives the first
equality in (F1T16.11). The finite Hurwitz map to the
\(b\)-pointed sphere has dimension \(b-3\); a period
image of dimension nine gives the second inequality.

The original rational primitive idempotent, its actual
integer-cleared Jacobian image, its field action and its
restricted theta form are retained. The induced maps on
the factor, parabolic cup pairing, Poincare biextension,
Hodge filtration and period differential are the identity
maps, with polarization multiplier one. The same identity
preserves the original Grothendieck and Adams realizations
in F01--F04. Thus the calculation excludes actual source
families for the prescribed sixfold corner; it introduces
no new curve, determinant component or algebraic detector.
